
\lesson{16/11}
[Weyl reflections]
\textbf{Def.} A root system $\Phi $ is the set of roots, i.e. the weight vectors
of the adjoint representation.\\
Its properties are
\begin{itemize}
    \item They span $\mathbb{R}^m$, where $m=\rank(\mathfrak{g})$;
    \item If $\alpha\in\Phi $, then $-\alpha \in\Phi $;
    \item If $\alpha,\beta \in\Phi$, $s_\alpha(\beta)\in\Phi$ (Weyl reflection);
    \item If $\alpha,\beta \in\Phi $, then $2\frac{\alpha\cdot\beta}{\alpha^2}\in
        \mathbb{Z}$.
\end{itemize}
\textbf{Example} $SU(3)$\\
The algebra of $SU(3)$ has $\vartheta =\frac{\pi}{3}$ and
$\rank(\mathfrak{su}(3))=2$; two out of the 8 generators are Cartan generators, so
in $\mathbb{R}^2$ this\footnote{what?} is represented as an \textbf{hexagon}.

Only $2$ roots are needed to generate $\mathbb{R}^2$ and are called \textit{
fundamental roots}.
\subsubsection{An ordering of the weight vectors}
Consider a representation $D$ of dimension $n$ with weight vectors
\[
    \mu = (\mu_1,...\mu_n):
\]
the ordering is defined as
\[
    \mu>\mu'\iff (\mu-\mu')_i >0,\;\text{with }i\text{ being the index of the 
    first nonzero component}.
\]
A weight vector is \textit{positive} if its first nonzero entry is positive.

In the case of $\mathfrak{su}(3)$, the vectors ``on the right'' of the hexagon
are positive, while the others are negative; the ``largest'' weight vector is
$(1,0)$, while the second largest is $\left(\frac{1}{2},\frac{\sqrt{3}}{2}\right)
$.

\textbf{Def.} A root is called \textit{simple} if it is positive and cannot be
written as the sum of two positive roots.

The number of simple roots is the same as the \textit{rank of the algebra}.

For $\mathfrak{su}(3)$, the two simple roots are
$\left(\frac{1}{2},\frac{\sqrt{3}}{2}\right)$ and
$\left(\frac{1}{2},-\frac{\sqrt{3}}{2}\right)$.

The difference of two simple roots $\alpha,\beta$ is never a root:
\[
    \alpha-\beta\not\in\Phi.
\]
It can be proven by contradiction, using $\alpha=(\alpha-\beta)+\beta$, that this
is not a \textit{simple root} and neither a root, and similarly for $\beta-
\alpha$. This means
\[
    E_{-\alpha}\ket{E_\beta}=0
\]
thus the number of times one can act with $E_\alpha^\dagger =E_{-\alpha}$ on
$\ket{E_\beta}$ is zero, which\footnote{unsure this is the reason} means $q=0$ (
in $\frac{\alpha\cdot \beta}{\alpha^2}=-\frac{1}{2}(p-q)$).
Similarly
\[
    E_{-\beta}\ket{E_\alpha}=0
\]
means $q'=0$ in $\frac{\beta\cdot \alpha}{\beta^2}=-\frac{1}{2}(p'-q')$. Taking
the product of these equations yields
\[
    \frac{(\alpha\cdot \beta)^2}{\beta^2\alpha^2}=\frac{pp'}{4}=\cos^2(
    \vartheta_{\alpha\beta})
\]
and its ratio
\[
    \frac{\beta^2}{\alpha^2}=\frac{p}{p'}
\]
Since $\vartheta_{\alpha\beta}>\frac{\pi}{2}$(??) from the first follows that
\[
    \cos(\vartheta_{\alpha\beta})=-\frac{1}{2}\sqrt{pp'}
\]

All positive roots can be written as
\[
    \phi = \sum_{\alpha}^{} k_\alpha \alpha,\;k_\alpha\in\mathbb{N}
\]
where $\alpha$ are simple roots. The quantity
\[
    k=\sum_{\alpha}^{} k_\alpha\geq 1
\] is called a [smth] index and its value is
\begin{itemize}
    \item $1$ only for simple roots;
    \item $>1$ for positive but not simple roots;
    \item $2$ if $\phi=\alpha+\beta$: since 
        \[
            2\frac{\alpha\cdot \beta}{\alpha^2}=-(p-q)=-p<0
        \] it depends on the algebra;
    \item Similarly for $3, >3$ up to $m=\rank(\mathfrak{g})$
\end{itemize}

\subsubsection{Cartan Matrix}
\textbf{Def.} A \textit{coroot} $\alpha^V$ of the root $\alpha$ is
\[
    \alpha^V=\frac{2\alpha}{\alpha^2}
\]
The Cartan Matrix $A\in\text{mat}(m\times m)$ is
\[
    A_{\alpha,\beta}=(\alpha,\beta^V)
\]
\textbf{Properties}
\begin{itemize}
    \item By definition, $A_{\alpha,\alpha}=2\;\forall\alpha$ root;
    \item For $\alpha\neq\beta$, $A_{\alpha\beta}\in\mathbb{Z}^{-}$ ($\neq0$) by
        construction;
    \item $\det(A)\neq0$;
    \item If $A_{\alpha,\beta}=0$, then $A_{\beta,\alpha}=0$, by construction;
    \item $\forall\alpha\neq\beta$, $A_{\alpha,\beta}\cdot A_{\beta,\alpha}\in
        \left\{0,1,2,3\right\}$: this means that the value of the entries cannot
        grow indefinitely; 
\end{itemize}
\subsection{Dynkin Diagram}
\begin{itemize}
    \item This is a diagram made with $m$ circles, each associated to one simple
        root;
    \item The circles can be connected with $A_{\alpha,\beta}\cdot
        A_{\beta,\alpha}\in\left\{0,1,2,3\right\}$ lines;
    \item If $A_{\alpha,\beta}\cdot A_{\beta,\alpha}=2,3$, the line
        has to be denoted by drawing an arrow \textit{entering the smallest root};
    \item Roots of the same length are coloured in the same way.
\end{itemize}
To each Cartan matrix correspond a Dynkin diagram e.g.
\[
    A=\begin{pmatrix}
        2&-1\\
        -1&2
    \end{pmatrix} 
\]
represents $\mathfrak{su}(3)$ and the diagram is on wikipedia.
\subsection{Classification of Dynkin diagrams}
There are only $4$ infinite families of diagrams with $m$ nodes that satisfy the
properties:
\begin{itemize}
    \item $A_m$ which are open, linear chains with single lines.
        It corresponds to $\mathfrak{su}(m+1)$;
    \item $B_m$ which are open, linear chains with a double line at
        the end, with the smallest root on the right extremum. It corresponds to
        $\mathfrak{so}(2m+1)$;
    \item $C_m$ which are open, linear chains with a double line at
        the end, with the smallest root on the second to last node.
        It corresponds to $\mathfrak{sp}(2m)$;
    \item $D_m$ which are open, linear chains with a ``triangle" at the right
        extremum. It corresponds to $\mathfrak{so}(2m)$.
\end{itemize}
There are also $5$ exceptional cases.
\begin{itemize}
    \item $G_2$ is the only diagram with $3$ connections between two nodes;
    \item $F_4$ is the $4$-nodes diagram with a double line in the middle;
    \item $E_6,E_7,E_8$ present a linear chain with a single deviating node
        at the third (from left) node.
\end{itemize}
These form the exceptional Lie Groups.

\subsection{Example: $SU(3)$}

